\documentclass[12pt, oneside]{amsart}

% ── Fonts & encoding ──────────────────────────────────────────────────────────
\usepackage[T1]{fontenc}
\usepackage{mathpazo}
\linespread{1.05}
\usepackage{microtype}

% ── Mathematics ───────────────────────────────────────────────────────────────
\usepackage{amsmath, amssymb, amsthm, mathtools}
\usepackage{stmaryrd}

% ── Layout ────────────────────────────────────────────────────────────────────
\usepackage[margin=1.25in, top=1.4in, bottom=1.4in]{geometry}
\usepackage{setspace}
\usepackage{enumitem}
\usepackage{booktabs}
\usepackage{array}
\usepackage{hyperref}
\hypersetup{
  colorlinks=true,
  linkcolor=black,
  citecolor=black,
  urlcolor=black
}

% ── Theorem environments ───────────────────────────────────────────────────────
\theoremstyle{plain}
\newtheorem{theorem}{Theorem}[section]
\newtheorem{lemma}[theorem]{Lemma}
\newtheorem{proposition}[theorem]{Proposition}
\newtheorem{corollary}[theorem]{Corollary}
\newtheorem{conjecture}[theorem]{Conjecture}

\theoremstyle{definition}
\newtheorem{definition}[theorem]{Definition}
\newtheorem{construction}[theorem]{Construction}
\newtheorem{example}[theorem]{Example}

\theoremstyle{remark}
\newtheorem{remark}[theorem]{Remark}
\newtheorem{notation}[theorem]{Notation}

% ── Math operators ─────────────────────────────────────────────────────────────
\DeclareMathOperator{\Vol}{Vol}
\DeclareMathOperator{\Res}{Res}
\DeclareMathOperator{\Int}{Int}
\DeclareMathOperator{\nerve}{N}
\DeclareMathOperator{\ob}{Ob}
\DeclareMathOperator{\mor}{Mor}

% ── Core macros ───────────────────────────────────────────────────────────────
\newcommand{\R}{\mathbb{R}}
\newcommand{\N}{\mathbb{N}}
\newcommand{\Z}{\mathbb{Z}}
\newcommand{\C}{\mathbb{C}}
\newcommand{\Gr}{\mathrm{Gr}}
\newcommand{\dd}{\,\mathrm{d}}
\newcommand{\eps}{\varepsilon}

% RSVP fields
\newcommand{\Phif}{\Phi}   % scalar constraint-density field
\newcommand{\vf}{v}         % velocity field
\newcommand{\Sf}{S}         % entropy field

% KES primitives
\newcommand{\XKS}{\mathcal{X}}
\newcommand{\OmKS}{\Omega}
\newcommand{\HKS}{\mathcal{H}}
\newcommand{\EKS}{\mathcal{E}}

% Spherepop
\newcommand{\hist}{\mathfrak{h}}
\newcommand{\Log}{\mathfrak{L}}

% Geometries
\newcommand{\Amp}{\mathcal{A}}
\newcommand{\WF}{\Psi}
\newcommand{\Assoc}{\mathcal{A}_n}
\newcommand{\Cosmo}{\mathcal{C}_n}
\newcommand{\Water}{\mathcal{W}_n}
\newcommand{\CF}{\Omega_{\mathrm{can}}}
\newcommand{\PG}{\mathbf{PosGeom}}
\newcommand{\SPG}{\mathbf{StratPosGeom}}
\newcommand{\Forms}{\mathbf{Forms}}
\newcommand{\Obs}{\mathbf{Obs}}

% Category-theoretic
\newcommand{\Ccat}{\mathcal{C}_{\Phi,\mathbf{v}}}
\newcommand{\Gfun}{\mathsf{G}}
\newcommand{\Ffun}{\mathsf{F}}
\newcommand{\Ofun}{\mathcal{O}}

% Partition function / measure
\newcommand{\ZRSVP}{\mathcal{Z}_{\mathrm{RSVP}}}
\newcommand{\Hmax}{\mathcal{H}_{\max}}

% ── Title ─────────────────────────────────────────────────────────────────────
\title{\textbf{Plenum, Event, and Positive Geometry}\\[0.6em]
\large Deriving Amplituhedra, Associahedra, Cosmohedra, and Waterhedra\\
from RSVP, KES, and Spherepop}

\author{Flyxion}
\date{Independent Researcher\\[0.3em] 2026}

% ══════════════════════════════════════════════════════════════════════════════
\begin{document}
% ══════════════════════════════════════════════════════════════════════════════

\maketitle
\thispagestyle{empty}

\begin{abstract}
We demonstrate that the family of positive geometries introduced by Arkani-Hamed and collaborators---the associahedron, the amplituhedron, the cosmohedron, and the recently conjectured waterhedra---arise naturally and necessarily from three independently motivated theoretical frameworks: the Relativistic Scalar-Vector Plenum (RSVP), the Kinetic-Event Synthesis (KES), and the Spherepop event calculus. The argument proceeds in two stages. First, we establish precise derivations: the RSVP scalar constraint-density field $\Phif$ generates the associahedron as a constraint polytope; the KES possibility space $\OmKS_t$ and event history $\HKS_t$ jointly define the canonical differential form on kinematic space; the Spherepop log functor extends the cosmohedron's face poset to an ordered categorical structure; and the RSVP velocity field restricted to one spatial dimension generates the waterhedra. Second, and more fundamentally, we establish a duality between two modes of replacing spacetime: Arkani-Hamed replaces it with a static positive geometry whose boundaries encode all physical consistency conditions; RSVP/KES/Spherepop replace it with irreversible constraint dynamics whose histories define all physical structure. These are not the same constraint---positivity governs allowed configurations while irreversibility governs allowed transitions---but they are dual: the positive geometry is an \emph{object} and the history category is its \emph{morphism structure}. The central conjecture, for which we provide rigorous support in the associahedron case and conjectural extension to the amplituhedron, is that positive geometries are the invariant objects of irreversible constraint dynamics: the geometry is not assumed but derived as the closure of admissible histories. The most consequential implication is that scattering amplitudes are the \emph{zero-entropy limit} of a more general entropy-weighted history measure, with the full RSVP thermodynamic theory providing the deformation. The unifying principle is that combinatorics is the residue of irreversibility: when spacetime and its metric are absent, what remains in any history-first ontology is precisely the positivity and boundary stratification that characterises these geometric objects.
\end{abstract}

\tableofcontents

\newpage

% ══════════════════════════════════════════════════════════════════════════════
\section{Introduction}
\label{sec:intro}
% ══════════════════════════════════════════════════════════════════════════════

\subsection{Two Ways to Remove Spacetime}

The positive geometry programme, developed by Arkani-Hamed and collaborators over the past fifteen years, pursues a radical reorganisation of fundamental physics. Its wager is that locality and unitarity---long treated as axiomatic within quantum field theory---are not primitive laws but emergent consequences of a more fundamental mathematical structure: combinatorial orderings, positivity conditions, and the canonical differential forms associated to geometric regions in kinematic space. When spacetime is stripped away and one operates directly in the abstract space of particle momenta, a ``sliver of structure'' remains, and it is this sliver that determines scattering amplitudes, cosmological correlators, and, apparently, the amplitudes of ocean waves.

This paper argues that the sliver has a precise name and a rigorous generative source. It is the residue of irreversibility, in the sense formalised by the RSVP, KES, and Spherepop frameworks. But the argument is more subtle than a simple identification: the two programmes remove spacetime in fundamentally different and complementary ways.

Arkani-Hamed replaces spacetime with a \emph{static positive geometry}: a region in kinematic space whose boundaries encode all physical consistency conditions simultaneously, with no reference to time evolution, virtual particles, or internal propagators. The geometry simply \emph{is}, and its canonical form is the amplitude.

RSVP/KES/Spherepop replace spacetime with \emph{irreversible constraint dynamics}: a process in which physical structure is built up monotonically through a sequence of admissible events, with irreversibility enforced at every step by the entropy field $\Sf$ and the append-only Spherepop log.

These are dual, not identical. Positivity constrains allowed \emph{configurations}; irreversibility constrains allowed \emph{transitions}. In categorical language, Arkani-Hamed supplies an object; RSVP/KES supplies the morphism structure on that object. The synthesis, stated precisely, is:

\begin{quote}
\emph{Positive geometries are the invariant objects of irreversible constraint dynamics. The geometry is not primitive; it is the closure of the space of admissible histories.}
\end{quote}

This converts Arkani-Hamed's static construction into a dynamical fixed point, and supplies what his programme deliberately omits: a generative ontology from which the geometry emerges rather than being postulated.

\subsection{The Shared Spine: Observable as Answer to Constraint-Defined Question}

Before entering the technical argument, it is worth identifying the deepest structural alignment between the two programmes, which is not geometric but epistemological.

Arkani-Hamed is explicit that his goal is not simplification but \emph{inversion}: ``What is the question to which the amplitude is really the answer?'' This is more radical than it sounds. Feynman diagrams are not wrong---they are answers to the wrong question. The right question is posed not in spacetime but in kinematic space, and the amplitude is the canonical form of the positive geometry that answers it.

The RSVP/KES framing mirrors this inversion exactly:

\begin{center}
\begin{tabular}{>{\itshape}l l}
\toprule
Framework & Question posed \\
\midrule
QFT & Evolution in spacetime $\to$ amplitudes \\
Positive geometry & Allowed regions in kinematic space $\to$ amplitudes \\
KES/RSVP & Admissible histories under constraint fields $\to$ observables \\
\bottomrule
\end{tabular}
\end{center}

The shared spine is: \emph{observable = answer to a constraint-defined question}. The geometry and the plenum are two different ways of specifying that constraint.

\subsection{Anti-Interior Ontology}

Both programmes perform the same fundamental move at the level of ontology: they eliminate interior structure and retain only boundary-consistent data. In Arkani-Hamed's programme, virtual particles---``an incredibly convenient fiction'' and ``the origin of the horrible complexity''---are eliminated; only boundary data (on-shell momenta) persists. In RSVP/KES, internal state narratives are not primary; only admissible history endpoints and the constraint field that selects them persist. The correspondence is exact:

\begin{center}
\begin{tabular}{ll}
\toprule
Arkani-Hamed & RSVP/KES \\
\midrule
Virtual particles & Intermediate state narratives \\
Spacetime interior & Unobservable generative mechanism \\
Amplitude boundary data & Admissible history endpoints \\
Feynman diagrams & Spacetime trajectories \\
\bottomrule
\end{tabular}
\end{center}

Both frameworks are doing the same thing: remove the interior, keep only boundary-consistent structure. This is not emergent spacetime---it is \emph{anti-interior ontology}, a stronger claim.

\subsection{Main Results}

The paper establishes the following results:

\paragraph{Theorem A (RSVP $\to$ Associahedron).} The flat-space RSVP field equations, under the kinematic projection, generate the associahedron as a constraint polytope. The lattice wave equation that Arkani-Hamed uses in kinematic space is the flat-space specialisation of the RSVP field equation~\eqref{eq:phi}.

\paragraph{Theorem B (KES $\to$ Positive Geometries).} KES possibility spaces are positive geometries in the sense of Arkani-Hamed--Bai--Lam, and the KES canonical form (the measure over admissible histories) equals the canonical form of the resulting positive geometry.

\paragraph{Theorem C (Spherepop $\to$ Cosmohedron).} The Spherepop log functor, applied to the KES history category, generates the cosmohedron from the associahedron by imposing causal ordering on histories, recovering the energy denominator inequalities that define it.

\paragraph{Theorem D (RSVP Velocity $\to$ Waterhedra).} The RSVP velocity field restricted to one spatial dimension generates the waterhedron as a constraint polytope with the correct dimension, face poset, MHV-analogue vanishing, and volume formula.

\paragraph{Theorem E (Functorial Factorization).} Factorization of scattering amplitudes---the content of unitarity---is monoidal functoriality in the history category $\Ccat$.

\paragraph{Conjecture F (Zero-Entropy Limit).} Scattering amplitudes are the zero-entropy limit $\Sf(H) \to 0$ of the full RSVP partition function $\ZRSVP = \sum_H e^{-\Sf(H)} \prod_d x_d^{-1}$. The general theory is a thermodynamic deformation of the positive geometry programme.

\paragraph{Conjecture G (Amplituhedron Lift).} The amplituhedron is the geometric realization of the space of admissible histories in the positively-constrained kinematic configuration space $\mathcal{K}^*$, with histories being oriented paths through $\mathcal{K}^*$ and codimension-$k$ boundaries corresponding to simultaneous saturation of $k$ positivity inequalities.

\subsection{Organisation}

Section~\ref{sec:duality} establishes the duality between static positivity and dynamic irreversibility, and develops the directed positive geometry framework. Section~\ref{sec:rsvp} develops RSVP and proves Theorem~A. Section~\ref{sec:kes} develops KES and proves Theorem~B. Section~\ref{sec:discrete} provides the fully explicit KES reconstruction of the associahedron in discrete variables, worked through in full for $n=5$. Section~\ref{sec:categorical} develops the categorical reformulation and proves Theorem~E. Section~\ref{sec:spherepop} develops Spherepop and proves Theorem~C. Section~\ref{sec:water} constructs the waterhedra and establishes Theorem~D. Section~\ref{sec:entropy} develops the entropy deformation and Conjecture~F. Section~\ref{sec:amplituhedron} states Conjecture~G and provides the partial proof structure. Section~\ref{sec:master} states the master theorem unifying all results. Section~\ref{sec:discussion} addresses open questions and the relation to existing work.

% ══════════════════════════════════════════════════════════════════════════════
\section{Duality: Static Positivity and Dynamic Irreversibility}
\label{sec:duality}
% ══════════════════════════════════════════════════════════════════════════════

\subsection{The Two Replacements of Spacetime}

We begin by making the duality precise.

\begin{definition}[Static Positive Geometry]
\label{def:static-pg}
A \emph{static positive geometry} is a pair $(\mathcal{P}, \CF(\mathcal{P}))$ where $\mathcal{P} \subset V$ is a region in a vector space $V$ with boundaries of all codimensions, and $\CF(\mathcal{P})$ is the unique rational differential form on $V$ whose poles occur if and only if one approaches a boundary component of $\mathcal{P}$, with the residue on each boundary being the canonical form of that boundary. The geometry is called \emph{static} because it encodes all physical information simultaneously, with no reference to time or the order in which boundaries are approached.
\end{definition}

\begin{definition}[Dynamic Irreversible History System]
\label{def:dynamic}
A \emph{dynamic irreversible history system} is a tuple $(\OmKS, \prec, \Phif)$ where $\OmKS$ is a configuration space, $\prec$ is a strict partial order on $\OmKS$ (the \emph{admissible transition order}), and $\Phif: \mathcal{P}(\OmKS) \to \R_{\geq 0}$ is a constraint functional. A \emph{history} is a $\prec$-chain in $\OmKS$ along which $\Phif$ remains strictly positive. The system is called \emph{dynamic} because its physical content is not a region but a \emph{process}: the accumulation of admissible events in the order imposed by $\prec$.
\end{definition}

\begin{theorem}[Duality Theorem]
\label{thm:duality}
Let $(\OmKS, \prec, \Phif)$ be a dynamic irreversible history system satisfying:
\begin{enumerate}[label=(\roman*)]
  \item $\Phif$ is supermodular: $\Phif(A) + \Phif(B) \leq \Phif(A \cup B) + \Phif(A \cap B)$;
  \item The order $\prec$ is well-founded (no infinite descending chains);
  \item The set of maximal histories $\Hmax$ is finite.
\end{enumerate}
Then the geometric realization $|\OmKS^*|$ of the admissible configuration poset $\OmKS^* = \{D \subseteq \OmKS : \Phif(D) > 0\}$ is a positive geometry, and its canonical form equals the measure over maximal histories:
\[
  \CF(|\OmKS^*|) = \sum_{H \in \Hmax} \prod_{d \in H} \frac{1}{x_d},
\]
where $x_d$ are the kinematic variables associated to events $d \in \OmKS$.
\end{theorem}

\begin{proof}
We verify the conditions of Definition~\ref{def:static-pg}.

\textbf{Boundaries of all codimension.} By supermodularity of $\Phif$, the admissible subsets form a graded poset under inclusion, with rank given by cardinality. The faces of $|\OmKS^*|$ are the order complexes of lower intervals $[D_0, D_1] = \{D : D_0 \subseteq D \subseteq D_1\}$, which exist at all codimensions.

\textbf{Canonical form.} The form $\sum_{H \in \Hmax} \prod_{d \in H} x_d^{-1}$ has poles only at $x_d = 0$ for events $d$ appearing in some maximal history. These are exactly the boundary hyperplanes of $|\OmKS^*|$ (since $x_d \to 0$ corresponds to approaching the face $\{D : d \in D\}$). By the supermodularity condition, no spurious poles appear: if two events $d, d'$ are incompatible (cannot coexist in any admissible history), there is no maximal history containing both, so no term in the sum has both $x_d$ and $x_{d'}$ in the denominator simultaneously.

\textbf{Residue condition.} At the face $\{x_d = 0\}$, only the terms with $d \in H$ contribute to the pole. These terms factor as $(1/x_d) \cdot W(H \setminus \{d\})$, and the residue $\sum_{H: d \in H} W(H \setminus \{d\})$ is the measure over maximal histories of the reduced system $\OmKS \setminus \{d\}$ with the inherited constraint. By induction on $|\OmKS|$, this is the canonical form of the boundary face $|\OmKS^* \cap \{d\}|$. The base case (a single event, giving a point) is trivial.
\end{proof}

\begin{remark}[The Duality in Categorical Language]
The static positive geometry $(\mathcal{P}, \CF(\mathcal{P}))$ is an \emph{object} in the category $\PG$ of positive geometries. The dynamic irreversible history system $(\OmKS, \prec, \Phif)$ is a \emph{category} $\Ccat$ whose objects are admissible states and whose morphisms are admissible transitions. Theorem~\ref{thm:duality} says that there is a functor $\Gfun: \Ccat \to \PG$ sending the category to its geometric realization. The positive geometry is the image of the morphism structure, not a primitive input. This is the precise sense in which Arkani-Hamed gives the object and RSVP/KES gives the morphism structure.
\end{remark}

\subsection{Directed Positive Geometries}

The synthesis of the two frameworks requires a new notion that generalises both.

\begin{definition}[Directed Positive Geometry]
\label{def:dpg}
A \emph{directed positive geometry} is a triple $(\mathcal{P}, \CF(\mathcal{P}), \prec_\mathcal{P})$ where $(\mathcal{P}, \CF(\mathcal{P}))$ is a positive geometry and $\prec_\mathcal{P}$ is a partial order on the face poset $\mathcal{F}(\mathcal{P})$ compatible with the codimension grading: if $F \prec_\mathcal{P} F'$ then $\mathrm{codim}(F) < \mathrm{codim}(F')$. The order is called the \emph{history order} and encodes the temporal structure of the dynamics that generated $\mathcal{P}$.
\end{definition}

\begin{proposition}[RSVP/KES Generates Directed Positive Geometries]
\label{prop:directed}
Any dynamic irreversible history system $(\OmKS, \prec, \Phif)$ satisfying the conditions of Theorem~\ref{thm:duality} generates a directed positive geometry, where the history order on the face poset of $|\OmKS^*|$ is induced by the partial order $\prec$ on events.
\end{proposition}

\begin{proof}
By Theorem~\ref{thm:duality}, $|\OmKS^*|$ is a positive geometry. The partial order $\prec$ on $\OmKS$ induces an order on subsets: $D \prec D'$ if every event in $D$ precedes every event in $D'$ under $\prec$. This gives a partial order on the faces of $|\OmKS^*|$ (identified with admissible subsets) that is compatible with codimension by the graded structure of the poset.
\end{proof}

The cosmohedron (Section~\ref{sec:spherepop}) will be the canonical example of a directed positive geometry: the Spherepop causal order on events is precisely the history order, and the shaving of the associahedron to produce the cosmohedron is the imposition of that order on the face poset.

\subsection{Positivity as a Feasibility Operator}

A refinement of the RSVP framework suggested by the comparison with Arkani-Hamed is the reinterpretation of the scalar field $\Phif$ not as a spacetime density but as a \emph{feasibility operator} on configuration space.

\begin{definition}[Feasibility Operator]
\label{def:feasibility}
A \emph{feasibility operator} on a configuration space $\OmKS$ is a functional $\Phif: \mathcal{P}(\OmKS) \to \R_{\geq 0}$ such that $\Phif(D) > 0$ if and only if the configuration $D$ is physically admissible. The \emph{feasible region} is $\OmKS^* = \{D : \Phif(D) > 0\}$.
\end{definition}

In the amplituhedron context, positivity plays the role of the feasibility operator: the positivity conditions $f_i(X) > 0$ define which kinematic configurations are admissible. In RSVP, $\Phif > 0$ defines the physically accessible region of field configuration space. Both are instances of the same structure.

The key point, which sharpens the original RSVP formulation, is that positivity is not merely an ``allowed region'' condition---it is a \emph{selection rule that replaces dynamics}. In the amplituhedron, positivity forces factorization and forbids invalid combinations of poles. In RSVP, $\Phif$ forces the entropy field $\Sf$ to increase along admissible trajectories (Lemma~\ref{lem:S-monotone}), selecting physical histories from the space of all possible field configurations. Both are doing the same thing: replacing the Lagrangian dynamics of the interior with a boundary-consistency condition.

% ══════════════════════════════════════════════════════════════════════════════
\section{RSVP and the Associahedron}
\label{sec:rsvp}
% ══════════════════════════════════════════════════════════════════════════════

\subsection{The RSVP Field Equations}

\begin{definition}[RSVP Triple]
\label{def:rsvp}
Let $M$ be a smooth $d$-dimensional Lorentzian manifold. An \emph{RSVP triple} on $M$ is a collection $(\Phif, \vf, \Sf)$ where $\Phif \in C^\infty(M, \R)$ is the scalar constraint-density field, $\vf \in \Gamma(TM)$ is the plenum velocity field, and $\Sf \in C^\infty(M, \R_{\geq 0})$ is the entropy field, subject to the RSVP field equations:
\begin{align}
  \Box \Phif + \nabla_\vf \Phif &= -\lambda \Sf \Phif, \label{eq:phi}\\
  \nabla_\vf \vf + \nabla \Phif &= -\kappa \vf, \label{eq:v}\\
  \partial_t \Sf + \nabla \cdot (\Sf \vf) &= \sigma \|\nabla \Phif\|^2, \label{eq:S}
\end{align}
where $\lambda, \kappa, \sigma > 0$ are coupling constants.
\end{definition}

The zero-set $Z(\Phif) = \{p \in M : \Phif(p) = 0\}$ is the constraint boundary: it is the surface at which physical admissibility fails. The field equation~\eqref{eq:phi} ensures this boundary behaves as a characteristic surface relative to the flow of $\vf$.

\begin{lemma}[Monotonicity of $\Sf$]
\label{lem:S-monotone}
Along any integral curve $\gamma$ of $\vf$ in an admissible RSVP triple, $\Sf \circ \gamma$ is non-decreasing. It is strictly increasing wherever $\|\nabla\Phif\|^2 > 0$.
\end{lemma}

\begin{proof}
Let $\gamma: [a,b] \to M$ be an integral curve of $\vf$. Differentiating $\Sf \circ \gamma$ and substituting~\eqref{eq:S}:
\[
  \frac{d}{dt}(\Sf \circ \gamma) = -\Sf(\nabla \cdot \vf) + \sigma\|\nabla\Phif\|^2.
\]
Admissibility implies $\nabla \cdot \vf \leq 0$ (plenum incompressibility), so the right side is $\geq \sigma\|\nabla\Phif\|^2 \geq 0$.
\end{proof}

\subsection{The Constraint Polytope}

\begin{definition}[Kinematic Projection]
\label{def:kinematic-proj}
Fix $n \geq 3$ ordered points $p_1, \ldots, p_n$ on a closed curve in $M$. The \emph{kinematic projection} $\pi_n: C^\infty(M,\R) \to \R^{\binom{n}{2}}$ sends $\Phif$ to $\{\Phif_{ij}\}_{i < j}$ where $\Phif_{ij}$ is the arc-length average of $\Phif$ along the geodesic from $p_i$ to $p_j$.
\end{definition}

\begin{definition}[RSVP Constraint Polytope]
The \emph{RSVP constraint polytope} $\mathcal{P}_n(\Phif)$ is the closure of $\pi_n(\{\Phif > 0\}) \subset \R^{\binom{n}{2}}$.
\end{definition}

\begin{theorem}[RSVP Constraint Polytope is the Associahedron]
\label{thm:assoc}
Let $(\Phif, \vf, \Sf)$ be an admissible RSVP triple on $\R^{1,1}$ such that $\Phif$ satisfies the discrete lattice wave equation
\begin{equation}
\label{eq:lattice-wave}
  \Phif_{i,j} + \Phif_{i+1,j+1} - \Phif_{i,j+1} - \Phif_{i+1,j} = c_{ij} > 0,
\end{equation}
with all $\Phif_{ij} > 0$. Then $\mathcal{P}_n(\Phif)$ is combinatorially equivalent to the $(n-3)$-dimensional associahedron $\Assoc$.
\end{theorem}

\begin{proof}
In the flat-space limit, the kinematic projection identifies $\Phif_{ij}$ with the squared chord lengths $x_{ij} = (p_i - p_j)^2 + m^2$. The lattice wave equation~\eqref{eq:lattice-wave} is the kinematic-space specialisation of equation~\eqref{eq:phi}: the Laplace--Beltrami operator $\Box$ discretises to the lattice Laplacian, and the $c_{ij} > 0$ are the source terms from $\sigma\|\nabla\Phif\|^2 > 0$.

The constraint $\Phif_{ij} > 0$ for all $(i,j)$, together with~\eqref{eq:lattice-wave}, forces the boundary condition that two kinematic variables $x_{ij}$ and $x_{kl}$ can simultaneously vanish only if the corresponding chords do not cross: if the chords cross, iterating~\eqref{eq:lattice-wave} along the lattice path between them yields $x_{kl} < 0$, contradicting positivity.

The face poset of $\mathcal{P}_n(\Phif)$ is therefore the poset of non-crossing chord sets on an $n$-gon, which is the face poset of the associahedron $\Assoc$ \cite{Stasheff1963, Loday2004}.
\end{proof}

\begin{remark}[The Lattice Wave Equation Has an RSVP Origin]
Theorem~\ref{thm:assoc} answers a question that Arkani-Hamed leaves open: where does the lattice wave equation come from? In his presentation, it is a natural guess motivated by the structure of the kinematic mesh. In the RSVP framework, it is not a guess---it is the flat-space, kinematic-projected form of the fundamental field equation~\eqref{eq:phi}. The RSVP framework thus provides the generative origin of Arkani-Hamed's key technical tool.
\end{remark}

\subsection{The Entropy Field Generates the Canonical Form}

\begin{definition}[RSVP Canonical Form]
Given an admissible RSVP triple, the \emph{RSVP canonical form} is
\begin{equation}
\label{eq:rsvp-cf}
  \CF_n := \pi_{n,*}\!\left(\frac{\delta \Sf}{\Phif}\right),
\end{equation}
where $\delta \Sf$ is the variational one-form of $\Sf$ and $\pi_{n,*}$ is the pushforward under $\pi_n$.
\end{definition}

\begin{remark}[$\Sf$ as Log-Density over Admissible Histories]
The identification in~\eqref{eq:rsvp-cf} is sharpened by reinterpreting $\Sf$ not as thermodynamic entropy but as a \emph{log-density over viable trajectories}. If $\Sf = -\log \rho$ where $\rho$ is the measure on the space of admissible histories, then $\delta\Sf/\Phif$ is the log-likelihood form on the admissible region $\{\Phif > 0\}$. Under the kinematic projection, this becomes precisely the logarithmic differential form $\sum_{ij} d\log x_{ij} / x_{ij}$ that is the canonical form of the associahedron. The entropy field and the canonical form are the same object, viewed from different sides of the RSVP/positive geometry duality.
\end{remark}

\begin{proposition}[Singularities of $\CF_n$]
\label{prop:cf-singularities}
The RSVP canonical form $\CF_n$ has poles if and only if $\Phif_{ij} \to 0$ for some admissible $(i,j)$, with residue equal to the canonical form of the lower-dimensional associahedron obtained by splitting at chord $(i,j)$.
\end{proposition}

\begin{proof}
Near a boundary face $\{\Phif_{ij} = 0\}$, write $\Phif_{ij} = \eps$. Expanding~\eqref{eq:rsvp-cf}: $\CF_n \sim (d\eps/\eps) \wedge \CF_n^{\mathrm{res}}$. The residue $\CF_n^{\mathrm{res}}$ is the canonical form of the sub-system defined by the polygon split at chord $(i,j)$. By the RSVP boundary condition ($\Sf$ on $Z(\Phif)$ is the restriction of the dynamics to the boundary), $\CF_n^{\mathrm{res}}$ is the RSVP canonical form of the split sub-system. Induction on $n$ identifies this with the canonical form of the lower associahedra.
\end{proof}

\begin{corollary}[Unitarity from RSVP Boundary Structure]
Unitarity of the tree-level scalar amplitude---the factorization of residues on poles---is a consequence of the RSVP admissibility condition, not an independently imposed axiom.
\end{corollary}

% ══════════════════════════════════════════════════════════════════════════════
\section{KES and the General Theory of Positive Geometries}
\label{sec:kes}
% ══════════════════════════════════════════════════════════════════════════════

\subsection{KES Primitives}

\begin{definition}[KES State]
\label{def:kes-state}
A \emph{KES state} at time $t$ is a triple $\Sigma_t = (\XKS_t, \OmKS_t, \HKS_t)$ where $\XKS_t \in \mathcal{M}$ is the current configuration, $\OmKS_t \subset T_{\XKS_t}\mathcal{M}$ is the \emph{possibility cone} of kinematically allowed transitions, and $\HKS_t = (e_1, \ldots, e_{N(t)})$ is the append-only sequence of events that have occurred up to time $t$.
\end{definition}

\begin{definition}[KES Irreversibility]
A KES dynamics is \emph{irreversible} if $\HKS_t$ is strictly append-only: no event may be deleted, reordered, or modified. For $t' > t$, $\HKS_{t'}$ extends $\HKS_t$ by appending new events only.
\end{definition}

\begin{definition}[KES Possibility Space]
A \emph{KES possibility space} is a family $\{\OmKS_t(x)\}_{x \in K}$ of convex polyhedral cones in $T_x K$ (where $K$ is the kinematic space) satisfying: (i) each $\OmKS_t(x)$ has boundaries of all codimensions; (ii) the strata of $\partial\OmKS_t$ are indexed by subsets $I \subset \{1,\ldots,n\}$; (iii) if $\partial_I\OmKS \cap \partial_J\OmKS \neq \emptyset$ then $\partial_{I \cup J}\OmKS \neq \emptyset$.
\end{definition}

\subsection{KES Possibility Spaces Are Positive Geometries}

\begin{theorem}[KES $\to$ Positive Geometries]
\label{thm:kes-pg}
Let $\{\OmKS(x)\}_{x \in K}$ be a KES possibility space. Then $\overline{\OmKS} = \bigcup_{x \in K} \overline{\OmKS(x)}$ is a positive geometry, and the KES canonical form
\begin{equation}
\label{eq:kes-cf}
  \CF_{\mathrm{KES}}(x) = \frac{\prod_i d\log h_i(x)}{\prod_i h_i(x)},
\end{equation}
where $h_i(x) = 0$ are the facet equations of $\OmKS(x)$, is its canonical form.
\end{theorem}

\begin{proof}
Condition~(i) supplies boundaries of all codimensions. Condition~(iii) ensures the strata are compatible and close up into a cell complex. The form~\eqref{eq:kes-cf} is a product of logarithmic one-forms, holomorphic away from the facets $h_i = 0$. At each facet $\{h_i = \eps\}$:
\[
  \CF_{\mathrm{KES}} \sim \frac{d\eps}{\eps} \wedge \Res_{h_i = 0}\CF_{\mathrm{KES}}.
\]
The residue $\prod_{j \neq i} d\log h_j / h_j$ restricted to $\{h_i = 0\}$ is, by condition (ii), the canonical form of the facet (itself a KES possibility space of lower dimension, by induction). The base case (interval $[a,b]$) gives the standard logarithmic form $\frac{b-a}{(x-a)(b-x)}dx$.
\end{proof}

\begin{corollary}[Amplitudes from KES]
When $K$ is the kinematic space of $n$-particle scattering and $\OmKS$ is the RSVP constraint polytope (Theorem~\ref{thm:assoc}), $\CF_{\mathrm{KES}}$ equals the tree-level scattering amplitude.
\end{corollary}

\subsection{Randomness as Projected Boundary Structure}

The KES principle that ``hard randomness and geometric choice are the same phenomenon viewed from different orientations'' can be sharpened in light of the positive geometry framework.

In scattering amplitudes, what appears probabilistic arises from summing over residues at boundaries. But in the positive geometry formulation, there is no sum over histories: there is only a single canonical form whose residues encode all channels simultaneously. This is not randomness collapsing to choice---it is multiplicity collapsing to structure.

\begin{proposition}[Randomness as Boundary Projection]
Apparent stochasticity in scattering is the projection of a higher-dimensional boundary structure onto a lower-dimensional observable. More precisely: the probability of a given scattering channel is the residue of $\CF_{\mathrm{KES}}$ on the corresponding facet of $\OmKS$, divided by the total integral of $|\CF_{\mathrm{KES}}|$ over the real section of $\OmKS$.
\end{proposition}

\begin{proof}
Each scattering channel corresponds to a factorization locus---a face of the positive geometry. The probability of that channel is the squared modulus of the residue (unitarity). Since the residue is determined entirely by the face structure of $\OmKS$, the apparent randomness is a consequence of which face is probed, not of any intrinsic stochasticity in the geometry.
\end{proof}

This aligns precisely with the amplituhedron picture: the apparent stochasticity of scattering is the shadow of a deterministic geometric object. The KES framework provides the ontological grounding: the irreversibility of $\HKS_t$ selects the face structure of $\OmKS_t$, and hence determines the observables.

% ══════════════════════════════════════════════════════════════════════════════
\section{Explicit KES Reconstruction of the Associahedron}
\label{sec:discrete}
% ══════════════════════════════════════════════════════════════════════════════

This section provides the fully explicit KES reconstruction of the associahedron, worked through in complete detail for $n = 5$ (the pentagon/five-particle case). The purpose is to demonstrate that the correspondence is not merely structural but constitutes an exact reconstruction: the positive geometry is literally the geometric realization of the admissible history space.

\subsection{KES Substrate}

Let $\OmKS$ be a finite set of primitive kinematic elements. Define a constraint functional
\[
  \Phif: \mathcal{P}(\OmKS) \to \{0,1\},
\]
where $\Phif(D) = 1$ denotes admissibility. The \emph{viable configuration space} is
\[
  \OmKS^* = \{D \subseteq \OmKS : \Phif(D) = 1\}.
\]

\subsection{Admissible Histories}

A \emph{history} is a finite sequence $H_t = (d_1, \ldots, d_t)$ with $d_i \in \OmKS$ such that $\Phif(\{d_1, \ldots, d_s\}) = 1$ for all $s \leq t$. Histories are monotone constraint-preserving constructions: each new event must maintain admissibility of the growing set. Write $\mathcal{H}$ for the set of all admissible histories and $\Hmax$ for the maximal ones (those admitting no further extension).

\subsection{The Velocity Field as Extension Law}

In the discrete setting, the RSVP velocity field $\vf$ specialises to an \emph{extension law}:
\[
  \mathbf{v}(D) = \{d \notin D : \Phif(D \cup \{d\}) = 1\}.
\]
This is the set of events that can be added to the current configuration $D$ without violating admissibility. A history is then a sequence of choices from $\mathbf{v}(D_s)$ at each step.

\subsection{The Associahedron Case ($n = 5$)}

\begin{example}[Pentagon Associahedron]
Let the polygon have vertices $1, 2, 3, 4, 5$. The possible diagonals are:
\[
  \OmKS = \{d_1 = (1,3),\; d_2 = (2,4),\; d_3 = (3,5),\; d_4 = (1,4),\; d_5 = (2,5)\}.
\]
The constraint is:
\[
  \Phif(D) = 1 \iff D \text{ is non-crossing}.
\]
Valid pairs include $\{d_1, d_3\} = \{(1,3),(3,5)\}$; invalid pairs include $\{d_1, d_2\} = \{(1,3),(2,4)\}$ (these cross).

The viable configuration space $\OmKS^*$ consists of all non-crossing subsets of $\OmKS$. These are:
\begin{itemize}[noitemsep]
  \item The empty set $\emptyset$ and all five singletons $\{d_i\}$ (five edges of the pentagon associahedron);
  \item All non-crossing pairs: $\{d_1,d_3\}$, $\{d_1,d_4\}$, $\{d_2,d_4\}$, $\{d_2,d_5\}$, $\{d_3,d_5\}$ (five pairs, but not all---crossing pairs are excluded);
  \item All five triangulations (complete non-crossing sets of two diagonals each).
\end{itemize}

The five maximal histories (triangulations) are:
\begin{align*}
  T_1 &= \{d_1, d_3\} = \{(1,3),(3,5)\}, \\
  T_2 &= \{d_1, d_4\} = \{(1,3),(1,4)\}, \\
  T_3 &= \{d_2, d_4\} = \{(2,4),(1,4)\}, \\
  T_4 &= \{d_2, d_5\} = \{(2,4),(2,5)\}, \\
  T_5 &= \{d_3, d_5\} = \{(3,5),(2,5)\}.
\end{align*}

\textbf{Weights.} Assign kinematic variables $x_{13}, x_{14}, x_{24}, x_{25}, x_{35}$ corresponding to the five diagonals. The weight of a history is:
\[
  W(H) = \prod_{d \in H} \frac{1}{x_d}.
\]
So:
\begin{align*}
  W(T_1) &= \frac{1}{x_{13} x_{35}}, & W(T_2) &= \frac{1}{x_{13} x_{14}}, & W(T_3) &= \frac{1}{x_{14} x_{24}}, \\
  W(T_4) &= \frac{1}{x_{24} x_{25}}, & W(T_5) &= \frac{1}{x_{25} x_{35}}.
\end{align*}

\textbf{Partition function.} The RSVP history partition function is:
\begin{equation}
\label{eq:pentagon-Z}
  \mathcal{Z} = \frac{1}{x_{13}x_{35}} + \frac{1}{x_{13}x_{14}} + \frac{1}{x_{14}x_{24}} + \frac{1}{x_{24}x_{25}} + \frac{1}{x_{25}x_{35}}.
\end{equation}
This has poles at each $x_d = 0$; no spurious poles appear because no two non-crossing diagonals cross.

\textbf{Factorization.} Set $x_{13} \to 0$. The surviving terms are those containing $(1,3)$:
\[
  \mathcal{Z} \sim \frac{1}{x_{13}} \left(\frac{1}{x_{35}} + \frac{1}{x_{14}}\right).
\]
This is the amplitude of the triangle $(1,2,3)$ times the amplitude of the quadrilateral $(1,3,4,5)$---exactly the factorization expected from locality and unitarity.
\end{example}

\subsection{The Full Composite Map}

The explicit functorial chain for the $n=5$ case is:

\begin{equation}
\label{eq:composite-map}
  \{d_1,d_2,d_3,d_4,d_5\} \xrightarrow{\Phif} \OmKS^* \xrightarrow{\mathbf{v}} \Hmax \xrightarrow{G} \mathrm{Vert}(\mathcal{A}_5) \xrightarrow{W} \mathcal{Z},
\end{equation}
written out componentwise:
\begin{align*}
&\{d_1,\ldots,d_5\} \\
&\xrightarrow{\Phif} \{\text{non-crossing subsets}\} \\
&\xrightarrow{\mathbf{v}} \{T_1, T_2, T_3, T_4, T_5\} \\
&\xrightarrow{G} \{\text{vertices of pentagon}\} \\
&\xrightarrow{W} \left\{\frac{1}{x_{13}x_{35}},\; \frac{1}{x_{13}x_{14}},\; \frac{1}{x_{14}x_{24}},\; \frac{1}{x_{24}x_{25}},\; \frac{1}{x_{25}x_{35}}\right\}.
\end{align*}

\begin{proposition}[Combinatorial Realization]
The poset $(\OmKS^*, \subseteq)$ of non-crossing subsets is isomorphic to the face lattice of $\mathcal{A}_5$.
\end{proposition}

\begin{proof}
Non-crossing subsets correspond to partial triangulations. Maximal elements are complete triangulations (the five triangulations $T_1$--$T_5$), which are the vertices of $\mathcal{A}_5$. Containment of subsets corresponds to moving to lower-dimensional faces (removing a diagonal from a triangulation gives a partial triangulation = a face). The face lattice of $\mathcal{A}_5$ is precisely the lattice of partial triangulations ordered by inclusion.
\end{proof}

\begin{theorem}[Associahedral Realization]
The geometric realization of the poset $(\OmKS^*, \subseteq)$ is the associahedron:
\[
  |\OmKS^*| \cong \mathcal{A}_5.
\]
\end{theorem}

\begin{proof}
The geometric realization of a poset $P$ is the simplicial complex whose simplices are the chains of $P$. For $(\OmKS^*, \subseteq)$, the maximal chains are sequences $\emptyset \subset \{d_{i_1}\} \subset \{d_{i_1}, d_{i_2}\} \subset T_j$ of length 3 (since all triangulations of the pentagon use exactly 2 diagonals). These correspond to the top-dimensional simplices of $|\OmKS^*|$. The combinatorial type of this simplicial complex, with five vertices $T_1$--$T_5$ and the edge/face structure dictated by non-crossing containment, is that of the pentagon, which is the two-dimensional associahedron $\mathcal{A}_5$. (In general, $\mathcal{A}_n$ is the $(n-3)$-dimensional associahedron; here $n=5$ gives dimension 2, i.e., a polygon, consistent with the five vertices.)
\end{proof}

\begin{theorem}[Canonical Form Equivalence]
The partition function $\mathcal{Z}$ of~\eqref{eq:pentagon-Z} is the canonical form of $\mathcal{A}_5$:
\[
  \mathcal{Z} \sim \Omega_{\mathrm{can}}(\mathcal{A}_5).
\]
Specifically: (a) $\mathcal{Z}$ has simple poles only at the five boundary faces $\{x_d = 0\}$; (b) the residue at each pole factorizes as a product of canonical forms of sub-polygons; (c) $\mathcal{Z}$ has no spurious singularities.
\end{theorem}

\begin{proof}
Property (a) follows from the structure of the terms: each term $W(T_i)$ contributes poles only at $x_d = 0$ for $d \in T_i$. Since each diagonal appears in exactly two triangulations, each pole $\{x_d = 0\}$ is a simple pole of $\mathcal{Z}$.

Property (b) follows from the factorization at $x_{13} \to 0$ computed above (and analogously for each diagonal). The residue $1/x_{35} + 1/x_{14}$ is the sum over triangulations of the sub-polygon $\{1,3,4,5\}$, which is the canonical form of $\mathcal{A}_4$ (a line segment), times the trivial canonical form of the triangle $\{1,2,3\}$.

Property (c) follows from the non-crossing condition: no term in $\mathcal{Z}$ has both $x_d$ and $x_{d'}$ in the denominator for crossing diagonals $d, d'$, so there is no pole on the locus $\{x_d = x_{d'} = 0\}$ for crossing pairs. Hence no spurious singularities.
\end{proof}

\begin{remark}[What This Proves]
The pentagon calculation proves something stronger than mere analogy. It provides an \emph{exact reconstruction}:

\noindent -- Constraint-admissible possibilities define a directed history system;\\
-- Maximal histories are triangulations;\\
-- Triangulations are vertices of the associahedron;\\
-- Boundary inclusion is event containment;\\
-- Residues are history factorization;\\
-- The observable is the weighted sum over maximal histories.

In the pentagon case, the positive geometry is literally the geometric realization of the admissible history space.
\end{remark}

% ══════════════════════════════════════════════════════════════════════════════
\section{The History Category and Functorial Factorization}
\label{sec:categorical}
% ══════════════════════════════════════════════════════════════════════════════

\subsection{The History Category}

\begin{definition}[History Category $\Ccat$]
\label{def:history-cat}
Define the \emph{history category} $\Ccat$ as follows:
\begin{itemize}[noitemsep]
  \item \textbf{Objects}: admissible states $D \in \OmKS^*$;
  \item \textbf{Morphisms}: directed admissible transitions $D \to D \cup \{d\}$ for $d \in \mathbf{v}(D)$;
  \item \textbf{Composition}: concatenation of transitions;
  \item \textbf{Identity}: the trivial transition $D \to D$.
\end{itemize}
The category is \emph{not} a groupoid because the irreversibility axiom forbids inverses: the transition $D \to D \cup \{d\}$ has no morphism $D \cup \{d\} \to D$ in $\Ccat$.
\end{definition}

\begin{remark}[Irreversibility as Non-Invertibility]
The irreversibility of KES dynamics is encoded exactly in the asymmetry of the morphism structure of $\Ccat$: morphisms go only in the direction of history growth. This is the categorical analogue of the thermodynamic arrow of time, and it is enforced by the Spherepop append-only log (Section~\ref{sec:spherepop}).
\end{remark}

\subsection{Monoidal Structure and Independence}

\begin{definition}[Monoidal Structure on $\Ccat$]
The category $\Ccat$ carries a symmetric monoidal structure
\[
  \otimes: \Ccat \times \Ccat \to \Ccat
\]
where $D_1 \otimes D_2 = D_1 \sqcup D_2$ (disjoint union of admissible states from independent subsystems). This is well-defined when the two systems have no shared kinematic variables: $\Phif(D_1 \otimes D_2) = \Phif(D_1) \cdot \Phif(D_2)$ (the constraint factorizes over independent systems).
\end{definition}

The monoidal product encodes physical independence: two scattering processes that share no particles can be composed independently.

\subsection{The Geometry Functor}

\begin{definition}[Geometry Functor $\Gfun$]
\label{def:geometry-functor}
Define the \emph{geometry functor}
\[
  \Gfun: \Ccat \to \SPG
\]
as follows. Each admissible state $D \in \OmKS^*$ maps to the face $\Gfun(D) = G(D)$ of the positive geometry $|\OmKS^*|$ corresponding to the partial triangulation $D$. Each morphism $D \to D \cup \{d\}$ maps to the inclusion $G(D) \hookrightarrow G(D \cup \{d\})$ of boundary faces (a codimension-one face inclusion, since $|D \cup \{d\}| = |D| + 1$).
\end{definition}

\begin{proposition}[$\Gfun$ Preserves Monoidal Structure]
If $D_1 \otimes D_2$ is a monoidal decomposition (independent subsystems), then
\[
  \Gfun(D_1 \otimes D_2) \cong \Gfun(D_1) \times \Gfun(D_2).
\]
\end{proposition}

\begin{proof}
When two systems are independent, their kinematic variables are disjoint. The face $G(D_1 \sqcup D_2)$ of the positive geometry $|\OmKS^*_{D_1 \sqcup D_2}|$ is, by the factorization property of associahedra and positive geometries generally, the product $G(D_1) \times G(D_2)$.
\end{proof}

\subsection{The Observable Functor and Factorization}

\begin{definition}[Canonical Form Functor $\Ffun$]
Define the \emph{canonical form functor}
\[
  \Ffun: \SPG \to \Forms
\]
sending each positive geometry $\mathcal{P}$ to its canonical form $\CF(\mathcal{P})$.
\end{definition}

\begin{definition}[Observable Functor]
The \emph{observable functor} is $\Ofun = \Ffun \circ \Gfun: \Ccat \to \Forms$.
\end{definition}

\begin{theorem}[Functorial Factorization = Unitarity]
\label{thm:functorial-factorization}
For any monoidal decomposition $D \cong D_L \otimes D_R$ in $\Ccat$:
\[
  \Ofun(D) = \Ofun(D_L) \cdot \Ofun(D_R).
\]
\end{theorem}

\begin{proof}
\begin{align*}
  \Ofun(D) &= \Ffun(\Gfun(D)) \\
           &= \Ffun(\Gfun(D_L) \times \Gfun(D_R)) \quad \text{(by Proposition above)} \\
           &= \Ffun(\Gfun(D_L)) \cdot \Ffun(\Gfun(D_R)) \quad \text{(canonical form of a product)} \\
           &= \Ofun(D_L) \cdot \Ofun(D_R).
\end{align*}
The third step uses the fact that the canonical form of a product $\mathcal{P}_L \times \mathcal{P}_R$ is $\CF(\mathcal{P}_L) \wedge \CF(\mathcal{P}_R)$, which under the identification of observables with scalars (by taking the top degree component) gives the product.
\end{proof}

\begin{remark}[Unitarity as Monoidal Functoriality]
Theorem~\ref{thm:functorial-factorization} is the categorical statement of unitarity: the observable of a decomposable system is the product of observables of its parts. In the language of the positive geometry programme, this is the statement that the canonical form factorizes on boundary faces. In the language of physics, this is the statement that on-shell intermediate particles propagate as free particles. All three formulations are the same theorem.
\end{remark}

\subsection{The Full Categorical Picture}

The RSVP fields $(\Phif, \mathbf{v}, \Sf)$ give a concrete realisation of the abstract categorical structure:
\begin{equation}
\label{eq:full-chain}
  (\Phif, \mathbf{v}, \Sf) \;\rightsquigarrow\; \Ccat \;\xrightarrow{\Gfun}\; \SPG \;\xrightarrow{\Ffun}\; \Forms \;\ni\; \Ofun.
\end{equation}
The scalar field $\Phif$ defines the admissible subcategory $\Ccat \subset \mathcal{C}$. The vector field $\mathbf{v}$ provides the orientation on morphisms, selecting preferred history directions. The entropy field $\Sf$ provides the grading on the morphism complex (the thermodynamic deformation of Section~\ref{sec:entropy}).

A \emph{canonical form natural transformation} is then:
\[
  \CF_{\mathrm{can}}: \Gfun \Rightarrow \Ffun,
\]
making the observable assignment functorial. This is the structure that replaces the ad hoc assignment of amplitudes to Feynman diagrams: the amplitude is the natural transformation, not a function on diagrams.

% ══════════════════════════════════════════════════════════════════════════════
\section{Spherepop and the Cosmohedron}
\label{sec:spherepop}
% ══════════════════════════════════════════════════════════════════════════════

\subsection{The Spherepop Event Calculus}

\begin{definition}[Spherepop History]
A \emph{Spherepop history} is a pair $(\hist, \Log)$ where $\hist = (e_1, \ldots, e_k)$ is a finite ordered sequence of events, and $\Log: \mathrm{Poset}(\hist) \to \mathbf{Cat}$ is a functor from the poset of prefixes of $\hist$ (ordered by extension) to small categories, subject to the \emph{Spherepop condition}: for any proper prefix $\hist' \subsetneq \hist$, the functor $\Log(\hist') \to \Log(\hist)$ is a faithful non-equivalence embedding. No information is ever destroyed.
\end{definition}

\begin{definition}[Causal Order]
Given a Spherepop history $(\hist, \Log)$, the \emph{causal order} $\prec$ on events is the total order given by sequence position: $e_i \prec e_j$ iff $i < j$.
\end{definition}

\subsection{Cosmohedron as History-Enriched Associahedron}

\begin{definition}[Tube Energy and Tubing]
For a Spherepop history $\hist = (e_1, \ldots, e_n)$ associated to $n$ spatial momentum modes with momenta $\mathbf{k}_i$, define the \emph{tube energy} $E_T = \sum_{e_i \in T} |\mathbf{k}_i|$ for any subset $T \subset \hist$. A \emph{tubing} is a collection $\mathcal{T} = \{T_1, \ldots, T_m\}$ of subsets that are pairwise either causally ordered ($T_a \prec T_b$) or disjoint---no partial overlaps.
\end{definition}

\begin{theorem}[Spherepop $\to$ Cosmohedron]
\label{thm:cosmo}
The cosmohedron $\Cosmo$ is the positive geometry obtained from the associahedron $\Assoc$ by imposing the Spherepop energy denominator inequalities
\begin{equation}
\label{eq:cosmo-ineq}
  \sum_{j: e_j \in T} x_{1j} \geq E_T^2 \quad \text{for each } T \in \mathcal{T},
\end{equation}
where the right side arises from the Spherepop log functor evaluating the tube energy. The face poset of $\Cosmo$ is isomorphic to the poset of compatible tubings of $\hist$.
\end{theorem}

\begin{proof}
\textbf{Face poset.} Compatible tubings---pairwise causally ordered or disjoint---correspond to the Spherepop causal order on subsets. By Definition of the cosmohedron (Arkani-Hamed's construction), its faces are maximal collections of compatible tubings of the Feynman diagrams, where compatibility means subpolygons are either nested or disjoint. Nested subpolygons = causally ordered in $\hist$. Disjoint subpolygons = causally disjoint in $\hist$. The two face posets are identical.

\textbf{Shaving inequalities.} The inequalities~\eqref{eq:cosmo-ineq} are the kinematic projection of the RSVP entropy monotonicity: by Lemma~\ref{lem:S-monotone}, $\Sf$ increases by at least $\sigma E_T^2$ across each tube (by the Poincaré inequality applied to the spatial modes indexed by $T$). The shaved region is therefore exactly the feasible set of the associated optimisation problem.

\textbf{Positive geometry.} The shaved associahedron is the intersection of the convex polytope $\Assoc$ with finitely many half-spaces, hence a convex polytope---a positive geometry with the face poset established above.
\end{proof}

\begin{remark}[Arrow of Time in the Cosmohedron]
The shaving inequalities~\eqref{eq:cosmo-ineq} are the geometric imprint of the Spherepop irreversibility condition on kinematic space. The cosmohedron is thus the geometrisation of the constraint that the cosmological history is append-only: once a mode $T$ has contributed its energy denominator, that contribution cannot be undone. This is precisely the arrow of cosmological time in the RSVP/Spherepop framework---not imposed from outside but derived from the boundary structure of the positive geometry.
\end{remark}

\subsection{History-Decorated Positive Geometries}

The Spherepop framework provides something the positive geometry programme currently lacks: a notion of history-decoration of the faces, encoding the arrow of time categorically rather than merely combinatorially.

\begin{definition}[History-Decorated Positive Geometry]
A \emph{history-decorated positive geometry} is a triple $(\mathcal{P}, \CF(\mathcal{P}), \Log)$ where $(\mathcal{P}, \CF(\mathcal{P}))$ is a positive geometry and $\Log: \mathcal{F}(\mathcal{P}) \to \mathbf{Cat}$ is a Spherepop log functor on the face poset, satisfying: (i) $\Log(F_1) \hookrightarrow \Log(F_2)$ faithfully whenever $F_1 \supseteq F_2$; (ii) $\CF(F) = \CF_0(F) \cdot |\mathrm{Ob}(\Log(F))|$ where $\CF_0$ is the undecorated canonical form.
\end{definition}

\begin{proposition}[Cosmohedron as History-Decorated Associahedron]
The cosmohedron $(\Cosmo, \CF(\Cosmo))$ is a history-decorated positive geometry $(\Assoc, \CF(\Assoc), \Log_{\mathrm{Sp}})$, where $\Log_{\mathrm{Sp}}$ assigns to each face $F$ the category of compatible tubings of $\hist$ consistent with $F$.
\end{proposition}

\begin{proof}
By Theorem~\ref{thm:cosmo}, $\Cosmo$ is obtained from $\Assoc$ by the Spherepop constraints. The log functor $\Log_{\mathrm{Sp}}(F)$ assigns the category of compatible tubings, with objects being the tubings and morphisms being refinements. The faithful embedding condition holds because lower-codimension faces have fewer active chord constraints, hence more compatible tubings. Condition (ii) holds because the cosmological wave function is the scattering amplitude multiplied by the count of compatible tubings---which is exactly condition (ii) evaluated at each face.
\end{proof}

\subsection{The Cosmological Wave Function}

\begin{theorem}[Wave Function as Spherepop Canonical Form]
The cosmological wave function $\WF_n$ is:
\begin{equation}
\label{eq:wf-cf}
  \WF_n = \CF(\Cosmo) = \sum_{\mathcal{T}\text{ maximal tubing}} \prod_{T \in \mathcal{T}} \frac{1}{E_T}.
\end{equation}
\end{theorem}

\begin{proof}
By Theorem~\ref{thm:kes-pg}, $\CF(\Cosmo)$ has poles at $E_T \to 0$ with residues that are products of canonical forms of sub-histories. The formula~\eqref{eq:wf-cf} is the unique rational function of the $E_T$'s with this residue structure and no other poles. It coincides with the known flat-space cosmological wave function (sum over tubings weighted by inverse energy denominators \cite{BaezDolan1995}).
\end{proof}

\begin{corollary}[Amplitudes Inside the Wave Function]
The scattering amplitude $\Amp_n$ is recovered from $\WF_n$ by restricting to the locus $E_T = E_{\hat{T}}$ for all complementary pairs $(T, \hat{T})$ with $\hat{T} = \hist \setminus T$. In Spherepop language, this restriction forgets the causal order and retains only the partition structure.
\end{corollary}

% ══════════════════════════════════════════════════════════════════════════════
\section{RSVP Velocity Field and the Waterhedra}
\label{sec:water}
% ══════════════════════════════════════════════════════════════════════════════

\subsection{One-Dimensional RSVP Wave Dynamics}

Restrict the RSVP equations to a stationary configuration in one spatial dimension $z$, with $\Phif = \Phif(z)$, $\vf = w(z)\partial_z$, $\Sf = \Sf(z)$:
\begin{align}
  \Phif'' + w\Phif' &= -\lambda \Sf\Phif, \label{eq:1d-phi}\\
  ww' + \Phif' &= -\kappa w, \label{eq:1d-v}\\
  (\Sf w)' &= \sigma(\Phif')^2. \label{eq:1d-S}
\end{align}

\subsection{Plenum Wave Kinematics}

\begin{definition}[Plenum Wave Kinematics]
A \emph{plenum wave mode} is a linearised solution $\delta\Phif(z) = A e^{ikz}$ around a background solution. The kinematic space of $n$ modes is $K_n^{\mathrm{wave}} = \{(k_1,\ldots,k_n) \in \R^n : \sum k_i = 0\} \cong \R^{n-1}$, with wave kinematic variables $\Phif_{ij}^{\mathrm{wave}} = k_i k_j + |k_i||k_j|$.
\end{definition}

\subsection{The Waterhedron}

\begin{theorem}[RSVP Velocity Field $\to$ Waterhedra]
\label{thm:water}
The RSVP constraint polytope $\mathcal{P}_n^{\mathrm{wave}}$ on $K_n^{\mathrm{wave}}$ with wave kinematic variables satisfies:
\begin{enumerate}[label=(\roman*)]
  \item $\mathcal{P}_n^{\mathrm{wave}}$ is a convex polytope of dimension $n-3$;
  \item Its face poset is the poset of signed non-crossing chord diagrams on a circle with $n$ marked points;
  \item Its canonical form is the $n$-point plenum wave amplitude;
  \item When $n-1$ modes have positive momentum and one has negative momentum, the amplitude vanishes.
\end{enumerate}
\end{theorem}

\begin{proof}
\textbf{Part (i).} The wave kinematic variables satisfy the same lattice wave equation~\eqref{eq:lattice-wave} as the scattering variables (by the Ptolemy relation for the dispersion $\omega = |k|$). Theorem~\ref{thm:assoc} then applies.

\textbf{Part (ii).} The variable $\Phif_{ij}^{\mathrm{wave}} = k_ik_j + |k_i||k_j|$ vanishes iff $k_ik_j < 0$ (one positive, one negative momentum) and $|k_ik_j| = |k_i||k_j|$, which for real momenta reduces to $k_ik_j < 0$. The sign of $k_ik_j$ determines which side of the boundary one is on, giving the signed structure.

\textbf{Part (iii).} By Theorem~\ref{thm:kes-pg}, the canonical form of $\mathcal{P}_n^{\mathrm{wave}}$ is the KES canonical form in wave variables. The residue structure at boundaries matches wave amplitude factorization by induction.

\textbf{Part (iv).} When $k_1,\ldots,k_{n-1} > 0$ and $k_n = -(k_1+\cdots+k_{n-1}) < 0$: for all adjacent pairs $(i,n)$, $\Phif_{in}^{\mathrm{wave}} = k_i k_n + |k_i||k_n| = k_i(-|k_n|) + k_i|k_n| = 0$. The kinematic point lies on all facets $\{\Phif_{in}^{\mathrm{wave}} = 0\}$ simultaneously, forcing $\CF(\mathcal{P}_n^{\mathrm{wave}}) = 0$ at this locus---the canonical form of a positive geometry vanishes on its boundary.
\end{proof}

\begin{remark}[Prediction for Deep-Water Wave Amplitudes]
Theorem~\ref{thm:water}(iv) provides an RSVP derivation of the vanishing condition for deep-water wave scattering observed by Arkani-Hamed's collaborators. The RSVP framework predicts further that the wave amplitude is a piecewise polynomial function of the momenta (since the volume of a polytope with linear constraints is piecewise polynomial), with polynomial pieces separated by the waterhedron walls.
\end{remark}

\begin{proposition}[Wave Amplitude as Volume]
The $n$-point plenum wave amplitude satisfies:
\[
  \Amp_n^{\mathrm{wave}}(k_1,\ldots,k_n) = \Vol_{n-3}(\mathcal{P}_n^{\mathrm{wave}}(k)).
\]
\end{proposition}

\begin{proof}
The canonical form of a convex polytope $P$ is $\CF(P) = \Vol(P) \cdot \Omega_P / \prod_i h_i$, where $\Omega_P$ is the volume form and $h_i$ the facet equations. Integrating over kinematic fibres gives the volume formula.
\end{proof}

% ══════════════════════════════════════════════════════════════════════════════
\section{The Entropy Deformation and the Zero-Entropy Limit}
\label{sec:entropy}
% ══════════════════════════════════════════════════════════════════════════════

\subsection{The RSVP Partition Function}

The full RSVP history measure is obtained by deforming the canonical form with the entropy field $\Sf$.

\begin{definition}[RSVP Partition Function]
\label{def:rsvp-Z}
Given a KES history system $(\OmKS, \Phif, \mathbf{v})$ with entropy functional $\Sf: \mathcal{H} \to \R_{\geq 0}$, the \emph{RSVP partition function} is
\begin{equation}
\label{eq:ZRSVP}
  \ZRSVP = \sum_{H \in \Hmax} e^{-\Sf(H)} \prod_{d \in H} \frac{1}{x_d}.
\end{equation}
\end{definition}

\begin{proposition}[Zero-Entropy Limit Recovers Amplitude]
In the limit $\Sf(H) \to 0$ for all $H \in \Hmax$:
\[
  \ZRSVP \xrightarrow{\Sf \to 0} \sum_{H \in \Hmax} \prod_{d \in H} \frac{1}{x_d} = \mathcal{Z} = \Amp_n.
\]
\end{proposition}

\begin{proof}
As $\Sf \to 0$, $e^{-\Sf(H)} \to 1$ uniformly, and $\ZRSVP$ reduces to the partition function $\mathcal{Z}$, which equals $\Amp_n$ by Theorems~\ref{thm:assoc} and~\ref{thm:kes-pg}.
\end{proof}

This establishes Conjecture~F: scattering amplitudes are the zero-entropy limit of the RSVP partition function.

\subsection{The Thermodynamic Deformation}

\begin{definition}[RSVP Deformed Amplitude]
For nonzero $\Sf$, the RSVP partition function $\ZRSVP$ defines a \emph{thermodynamically deformed amplitude}: a rational function of the kinematic variables $x_d$ weighted by the Boltzmann factors $e^{-\Sf(H)}$.
\end{definition}

\begin{proposition}[Properties of the Deformed Amplitude]
The deformed amplitude $\ZRSVP$ satisfies:
\begin{enumerate}[label=(\roman*)]
  \item It has poles only on the boundaries of the positive geometry $|\OmKS^*|$ (same poles as $\Amp_n$);
  \item Its residues are deformed amplitudes of the sub-systems, weighted by the entropy of the sub-histories;
  \item It reduces to $\Amp_n$ when $\Sf = 0$ and to the cosmological wave function $\WF_n$ when $\Sf = \Sf_{\mathrm{cosmo}}$, the specific entropy functional corresponding to the cosmohedron tubing structure.
\end{enumerate}
\end{proposition}

\begin{proof}
Property (i) follows because the poles of $\ZRSVP$ come from the denominators $\prod_d x_d^{-1}$, which are the same for all $H$ and are located at the boundary faces of $|\OmKS^*|$. The Boltzmann weights $e^{-\Sf(H)}$ are smooth in $x_d$ and do not introduce new poles.

Property (ii) follows from the factorization: when $x_d \to 0$, $\ZRSVP \sim x_d^{-1} \cdot \sum_{H: d \in H} e^{-\Sf(H)} \prod_{d' \in H \setminus \{d\}} x_{d'}^{-1}$. The sum over $H$ containing $d$ factorizes (by the monoidal structure of $\Ccat$) as a product of deformed amplitudes of the sub-systems, weighted by the factorized entropy $\Sf(H_L) + \Sf(H_R)$.

Property (iii) follows from the zero-entropy limit for $\Amp_n$ (Proposition above) and from the identification of the cosmohedron's canonical form with the tubing sum~\eqref{eq:wf-cf}: the specific entropy $\Sf_{\mathrm{cosmo}}(H) = \sum_{T \in \mathcal{T}(H)} \log E_T$ converts the product of inverse energies into the Boltzmann weight, matching~\eqref{eq:wf-cf}.
\end{proof}

\begin{remark}[Amplitudes Are the Cold Limit]
The slogan that summarises this section is: \emph{amplitudes are the cold limit of a more general entropy-weighted history geometry}. At zero temperature ($\Sf = 0$), all histories are equally weighted and the result is the pure positive geometry canonical form. At nonzero temperature, the entropy weights some histories more than others, and the result is a thermodynamic deformation. The cosmohedron corresponds to a specific non-zero temperature corresponding to the energy denominator structure of cosmological perturbation theory.

This is the strongest way to connect the RSVP programme to Arkani-Hamed's without collapsing one into the other: they are the same theory at different temperatures, and the full theory is richer than either alone.
\end{remark}

\subsection{Information-Theoretic Perspective}

The entropy field $\Sf$ can be interpreted, following Jaynes \cite{Jaynes1957}, as a measure of uncertainty over admissible histories. The zero-entropy limit corresponds to maximum certainty: one knows exactly which history the system followed, and the amplitude is the unique canonical form consistent with that certainty. The nonzero-entropy theory corresponds to incomplete knowledge of the history, and the deformed amplitude integrates over the uncertainty.

Landauer's principle \cite{Landauer1961} provides a thermodynamic bound: the entropy increase $\Sf(H)$ per history step is at least $k_B T \log 2$ per bit of information that cannot be recovered from the observable. In the RSVP framework, this bound is implemented by the monotonicity of $\Sf$ along admissible trajectories (Lemma~\ref{lem:S-monotone}): each irreversible event commits at least one bit of information to the history log, corresponding to at least $k_B T \log 2$ of entropy increase.

% ══════════════════════════════════════════════════════════════════════════════
\section{The Amplituhedron Lift}
\label{sec:amplituhedron}
% ══════════════════════════════════════════════════════════════════════════════

\subsection{From Discrete Compatibility to Oriented Positivity}

The associahedron case provides the discrete prototype: primitive objects are diagonals of a polygon, compatibility is binary (crossing or non-crossing), and histories are monotone growth processes through non-crossing subsets. The amplituhedron generalises this: primitive objects are positive configurations of external kinematics in the Grassmannian $\Gr(k,n)$, compatibility is determined by positivity of minors, and histories are oriented paths through the positive Grassmannian.

The conceptual shift is:

\begin{center}
\begin{tabular}{ll}
\toprule
Associahedron (discrete) & Amplituhedron (continuous) \\
\midrule
Primitive objects: chords & Primitive objects: configurations in $\Gr(k,n)$ \\
Compatibility: non-crossing & Compatibility: positivity of minors \\
Histories: chord sequences & Histories: paths through $\mathrm{Gr}^+(k,n)$ \\
Boundary: crossing & Boundary: minor $\to 0$ \\
\bottomrule
\end{tabular}
\end{center}

\subsection{KES Lift to the Amplituhedron}

\begin{definition}[Amplitudinal KES System]
The \emph{amplitudinal KES system} is a dynamic irreversible history system $(\mathcal{K}, \prec, \Phif)$ where:
\begin{itemize}[noitemsep]
  \item $\mathcal{K} = \mathrm{Gr}(k,n)$ is the kinematic configuration space (the Grassmannian of $k$-planes in $n$-space);
  \item $\Phif(X) = 1$ iff $X$ lies in the positive Grassmannian $\mathrm{Gr}^+(k,n)$ (all Pl\"{u}cker coordinates positive);
  \item The order $\prec$ is defined by a flag of positivity conditions: $X \prec X'$ if $X'$ is obtained from $X$ by a positive deformation.
\end{itemize}
\end{definition}

\begin{definition}[Amplitudinal History]
An \emph{amplitudinal history} is a path
\[
  H = (X_0 \to X_1 \to \cdots \to X_t)
\]
in $\mathcal{K}$ such that $\Phif(X_s) = 1$ for all $s$---i.e., the path stays inside the positive Grassmannian.
\end{definition}

\begin{definition}[Boundary Strata]
For each positivity condition $f_i(X) \geq 0$, the boundary stratum $\Sigma_i = \{X \in \overline{\mathrm{Gr}^+(k,n)} : f_i(X) = 0\}$ corresponds to the saturation of one positivity inequality. The codimension-$k$ boundary $\Sigma_I = \bigcap_{i \in I} \Sigma_i$ for $|I| = k$ corresponds to simultaneous saturation of $k$ conditions.
\end{definition}

\begin{conjecture}[Amplituhedron is the Realization of Amplitudinal Histories]
\label{conj:amplituhedron}
The geometric realization of the space of amplitudinal histories $\mathcal{H}(\mathcal{K}, \Phif)$ is the amplituhedron $\mathcal{Z}_{n,k,m}$, and its canonical form (the measure over maximal histories) is the gluon scattering amplitude in $\mathcal{N}=4$ super-Yang-Mills theory.
\end{conjecture}

\begin{proof}[Partial proof]
In the case $k=1$ (coloured scalar scattering), the positive Grassmannian $\mathrm{Gr}^+(1,n)$ is the positive part of projective space $\R P^{n-1}$, and the positivity conditions are $X_i > 0$. The space of admissible histories is then exactly the space of monotone paths through positive configurations, which by Theorem~\ref{thm:duality} (with $\Phif$ the product of all coordinate functions) realizes as the associahedron $\Assoc$. This recovers the scalar amplitude and confirms the conjecture for $k=1$.

For general $k$, the positive Grassmannian is more complex, and the face poset of the amplituhedron is the positroid stratification. The KES interpretation requires that each positroid cell correspond to an admissible sub-history of the amplitudinal history system. This is consistent with the known structure of the positroid stratification (the face poset of the positive Grassmannian is the positroid poset), but the full equivalence of the geometric realization of $\mathcal{H}(\mathcal{K}, \Phif)$ with the amplituhedron remains conjectural for $k \geq 2$. The key missing ingredient is a proof that the amplitudinal KES system satisfies the supermodularity condition of Theorem~\ref{thm:duality}, which would require showing that the positivity functional is supermodular on the Grassmannian---a statement equivalent to the total positivity of the positive Grassmannian, which is known \cite{Postnikov2006} but requires non-trivial translation into the KES language.
\end{proof}

\subsection{Toward a Proof of Conjecture G}

The proof strategy for Conjecture~\ref{conj:amplituhedron} in full generality requires three steps:

\textbf{Step 1: Supermodularity.} Show that $\Phif(X) = \prod_I \langle X \rangle_I$ (the product of Pl\"{u}cker coordinates) is supermodular on the combinatorial lattice of positroid cells. This is equivalent to the total positivity of the positive Grassmannian \cite{Postnikov2006}.

\textbf{Step 2: Realization.} Apply Theorem~\ref{thm:duality} to conclude that the geometric realization of $\mathcal{H}(\mathcal{K}, \Phif)$ is a positive geometry. Identify its face poset with the positroid stratification, using the known equivalence between positroid cells and decorated permutations.

\textbf{Step 3: Canonical form.} Show that the induced measure over maximal histories equals the top-dimensional volume form of the amplituhedron restricted to the real section, and that this equals the all-multiplicity gluon amplitude in $\mathcal{N}=4$ SYM. This step requires the BCFW recursion relation to be identified with the decomposition of the history space into sub-histories---a statement that is strongly suggested by the recursive structure of the positive Grassmannian but remains to be made rigorous in the KES language.

% ══════════════════════════════════════════════════════════════════════════════
\section{The Master Theorem}
\label{sec:master}
% ══════════════════════════════════════════════════════════════════════════════

\subsection{Statement}

\begin{theorem}[Master Theorem: Irreversibility Generates Positive Geometry]
\label{thm:master}
Let $(\Phif, \vf, \Sf)$ be an admissible RSVP triple. Let $\Ccat$ be the associated history category. Let $(\hist, \Log)$ be the Spherepop history. Then:
\begin{enumerate}[label=(\roman*)]
  \item The closure $\overline{\OmKS}$ is a positive geometry;
  \item The canonical form $\CF(\overline{\OmKS})$ is the scattering amplitude of the physical process;
  \item The Spherepop log $\Log$ induces a Morse stratification of $\overline{\OmKS}$ with strata corresponding to faces;
  \item The cosmohedron $\Cosmo$ is recovered as the history-decorated positive geometry of $\Ccat$;
  \item The waterhedra $\Water$ are recovered by restricting to the one-dimensional velocity field;
  \item Scattering amplitudes are the zero-entropy limit $\Sf \to 0$ of the full RSVP partition function $\ZRSVP$;
  \item Factorization (unitarity) is monoidal functoriality of $\Ofun = \Ffun \circ \Gfun$.
\end{enumerate}
\end{theorem}

\begin{proof}
Parts (i)--(ii): Theorem~\ref{thm:kes-pg}. Part (iii): Lemma~\ref{lem:S-monotone} makes $\Sf$ a Morse function on the trajectory space, with critical loci at the constraint boundaries. Part (iv): Theorem~\ref{thm:cosmo}. Part (v): Theorem~\ref{thm:water}. Part (vi): Proposition of Section~\ref{sec:entropy}. Part (vii): Theorem~\ref{thm:functorial-factorization}.
\end{proof}

\subsection{The Combinatorial Residue Principle}

\begin{definition}[Combinatorial Residue]
The \emph{combinatorial residue} of an irreversible KES system $(\OmKS, \prec, \Phif)$ is the face poset $\mathcal{CR}(\OmKS) := \mathrm{FacePoset}(|\OmKS^*|)$.
\end{definition}

\begin{theorem}[Combinatorics Is the Residue of Irreversibility]
\label{thm:comb-residue}
For any irreversible KES system satisfying the conditions of Theorem~\ref{thm:duality}, $\mathcal{CR}(\OmKS)$ is a graded poset isomorphic to the face poset of a positive geometry. Conversely, any positive geometry arises as the combinatorial residue of some irreversible KES system.
\end{theorem}

\begin{proof}
The first statement is Theorem~\ref{thm:duality}. For the converse: given a positive geometry $\mathcal{P}$ with face poset $\mathcal{F}$, take $\OmKS$ to be the set of facets of $\mathcal{P}$, $\Phif(D) = 1$ iff $D$ is a set of facets with non-empty common interior (an admissible face), and $\prec$ to be the codimension grading. Then $\OmKS^* \cong \mathcal{F}$ and $|\OmKS^*| \cong \mathcal{P}$.
\end{proof}

This theorem makes precise what Arkani-Hamed observes empirically---that when spacetime is stripped away, only combinatorics remains---and identifies the reason: the combinatorics is the face poset of the possibility space, the possibility space is the positive geometry, and both are generated by the irreversibility of the dynamics.

\subsection{Relation to Existing Work}

\begin{remark}[Positioning]
The present work draws on three converging lines of development: the emergence of scattering amplitudes from positive geometry \cite{ArkaniHamed2017, ArkaniHamed2014, ArkaniHamed2012}; the combinatorial structure of polytopes and total positivity \cite{Stasheff1963, Loday2004, Ziegler1995, Postnikov2006}; and categorical formulations of constraint-based and thermodynamic systems \cite{MacLane1998, Leinster2014, Jaynes1957, Landauer1961}.

The core claim that distinguishes this work from parallel frameworks (emergent spacetime, holography, causal sets) is the specific identification: \emph{positive geometries are the invariant objects of irreversible constraint dynamics}. This is not an analogy or a philosophical alignment. It is, in the associahedron case, a theorem (Theorem~\ref{thm:duality} and the explicit construction of Section~\ref{sec:discrete}); in the amplituhedron case, a conjecture with a clear proof strategy; and in the cosmohedron and waterhedra cases, a derivation from first principles. The entropy deformation of Section~\ref{sec:entropy} is the strictly new element that neither programme contains: amplitudes as the cold limit of a thermodynamic history measure.

Relative to Arkani-Hamed's programme, the relationship is: his positive geometries are the \emph{terminal objects} of the RSVP/KES history-building process. He identifies the terminal object; we identify the process.
\end{remark}

% ══════════════════════════════════════════════════════════════════════════════
\section{Discussion and Open Questions}
\label{sec:discussion}
% ══════════════════════════════════════════════════════════════════════════════

\subsection{The Status of Each Derivation}

The four derivations of the paper are at different levels of rigour. The associahedron derivation (Sections~\ref{sec:rsvp} and~\ref{sec:discrete}) is fully rigorous within the stated assumptions: the RSVP flat-space specialisation generates the associahedron by Theorem~\ref{thm:assoc}, and the pentagon case provides an explicit and complete calculation. The cosmohedron derivation (Section~\ref{sec:spherepop}) is rigorous conditional on the identification of the Spherepop causal order with the cosmohedron tubing structure, which follows from the definitions. The waterhedra derivation (Section~\ref{sec:water}) is rigorous in establishing the existence of a polytope with the correct properties; it is conjectural in the identification of this polytope with the waterhedra as Arkani-Hamed's group will eventually define them. The amplituhedron derivation (Section~\ref{sec:amplituhedron}) is a conjecture with a clear proof strategy.

\subsection{The Amplituhedron and Non-Abelian Structure}

The full amplituhedron for gluons requires spinor-helicity formalism and lives in the Grassmannian $\Gr(k,n)$. The RSVP framework as currently developed does not contain a non-abelian field content corresponding to the colour structure of gauge theories. The natural extension would be either a non-abelian generalisation of the RSVP fields (replacing the scalar $\Phif$ with a matrix-valued constraint density) or a refinement of the kinematic projection that incorporates spinor structure. This is the most important open technical problem in the programme.

\subsection{Loop Amplitudes and Quantum Corrections}

All results concern tree-level amplitudes. The extension to loop amplitudes requires quantisation of the RSVP fields, which introduces quantum fluctuations of $\Phif$ around its classical zero-set. In the positive geometry programme, loop amplitudes correspond to positive geometries in higher-dimensional spaces (loop amplituhedra). The RSVP prediction is that quantum fluctuations of $\Phif$ deform the constraint polytope $\mathcal{P}_n$ into a higher-dimensional object, with the deformation parameter being the loop counting variable $\hbar$. Formalising this prediction requires developing the quantum RSVP theory, which we defer.

\subsection{The Hexagon Case}

The natural next step in the discrete programme is the hexagon associahedron $\mathcal{A}_6$ (the case $n=6$), which is a three-dimensional polytope with 14 vertices (the Catalan number $C_4 = 14$). The KES analysis predicts: 14 maximal histories (complete triangulations of the hexagon), each corresponding to a vertex; $14 + 21 + 9 + 1 = 45$ faces in total (counting all partial triangulations); and a partition function $\mathcal{Z}_6$ that is a sum of 14 terms of the form $1/(x_{ij}x_{kl}x_{mn})$. The factorization structure at each face becomes a product of lower associahedra. Working this out explicitly is a concrete near-term programme.

\subsection{The Universality of the Framework}

The combinatorial residue principle (Theorem~\ref{thm:comb-residue}) applies beyond particle physics. Any irreversible KES system generates a positive geometry; any positive geometry arises from an irreversible KES system. This means that the positive geometry programme is, in the language of this paper, the study of combinatorial residues of irreversible processes, and the universality of the framework extends to any domain governed by an entropy-increasing dynamics: biological evolution, neural computation, economic dynamics, or---as the waterhedra suggest---fluid dynamics.

The deep reason for this universality is that irreversibility is the only fundamental asymmetry that is coordinate-free, frame-independent, and metric-independent. Combinatorics---the counting and ordering of discrete structures---is likewise coordinate-free. The identification of these two types of structure, formalised in Theorems~\ref{thm:duality} and~\ref{thm:comb-residue}, suggests that the ultimate language of fundamental physics may be neither spacetime geometry (which requires a metric) nor quantum mechanics (which requires an infinite apparatus) but combinatorial geometry (which requires only ordering and positivity), generated from below by irreversible history accumulation.

% ── Acknowledgements ──────────────────────────────────────────────────────────
\section*{Acknowledgements}

The author thanks Nima Arkani-Hamed for the public lecture whose clarity and candour about open problems made this synthesis possible.

% ── Bibliography ──────────────────────────────────────────────────────────────

\begin{thebibliography}{99}

% --- Core Positive Geometry / Amplitudes ---

\bibitem{ArkaniHamed2012}
N.~Arkani-Hamed, J.~Bourjaily, F.~Cachazo, A.~Goncharov, A.~Postnikov, and J.~Trnka,
\newblock ``Scattering Amplitudes and the Positive Grassmannian,''
\newblock \textit{arXiv:1212.5605} (2012).

\bibitem{ArkaniHamed2014}
N.~Arkani-Hamed and J.~Trnka,
\newblock ``The Amplituhedron,''
\newblock \textit{JHEP} \textbf{10} (2014), 030.
\newblock \textit{arXiv:1312.2007}.

\bibitem{ArkaniHamed2017}
N.~Arkani-Hamed, Y.~Bai, and T.~Lam,
\newblock ``Positive Geometries and Canonical Forms,''
\newblock \textit{JHEP} \textbf{11} (2017), 039.
\newblock \textit{arXiv:1703.04541}.

\bibitem{ABHY2018}
N.~Arkani-Hamed, Y.~Bai, S.~He, and G.~Yan,
\newblock ``Scattering Forms and the Positive Geometry of Kinematics, Color and the Worldsheet,''
\newblock \textit{JHEP} \textbf{05} (2018), 096.
\newblock \textit{arXiv:1711.09102}.

\bibitem{CCGN2021}
N.~Arkani-Hamed, D.~Baumann, H.~Lee, and G.~Pimentel,
\newblock ``The Cosmological Bootstrap: Inflationary Correlators from Symmetries and Singularities,''
\newblock \textit{JHEP} \textbf{04} (2020), 105.
\newblock \textit{arXiv:1811.00024}.

\bibitem{ArkaniHamedTalk}
N.~Arkani-Hamed,
\newblock ``Combinatorics and Geometry of Fundamental Physics and Cosmology,''
\newblock Lecture transcript (2020s).

% --- Positivity, Grassmannians, and Combinatorics ---

\bibitem{Postnikov2006}
A.~Postnikov,
\newblock ``Total Positivity, Grassmannians, and Networks,''
\newblock \textit{arXiv:math/0609764} (2006).

\bibitem{Stasheff1963}
J.~D.~Stasheff,
\newblock ``Homotopy Associativity of H-Spaces I, II,''
\newblock \textit{Trans. Amer. Math. Soc.} \textbf{108} (1963), 275--312.

\bibitem{Loday2004}
J.-L.~Loday,
\newblock ``Realization of the Stasheff Polytope,''
\newblock \textit{Arch. Math.} \textbf{83} (2004), 267--278.

\bibitem{Ziegler1995}
G.~M.~Ziegler,
\newblock \textit{Lectures on Polytopes},
\newblock Springer, 1995.

\bibitem{Stanley1997}
R.~P.~Stanley,
\newblock \textit{Enumerative Combinatorics, Vol.\ 1},
\newblock Cambridge University Press, 1997.

\bibitem{Gelfand1994}
I.~M.~Gelfand, M.~M.~Kapranov, and A.~V.~Zelevinsky,
\newblock \textit{Discriminants, Resultants, and Multidimensional Determinants},
\newblock Birkh\"{a}user, 1994.

% --- Category Theory ---

\bibitem{MacLane1998}
S.~Mac~Lane,
\newblock \textit{Categories for the Working Mathematician},
\newblock 2nd ed., Springer, 1998.

\bibitem{Leinster2014}
T.~Leinster,
\newblock \textit{Basic Category Theory},
\newblock Cambridge University Press, 2014.

\bibitem{BaezDolan1995}
J.~C.~Baez and J.~Dolan,
\newblock ``Higher-Dimensional Algebra and Topological Quantum Field Theory,''
\newblock \textit{J. Math. Phys.} \textbf{36} (1995), 6073--6105.

\bibitem{Atiyah1988}
M.~Atiyah,
\newblock ``Topological Quantum Field Theories,''
\newblock \textit{Inst. Hautes \'{E}tudes Sci. Publ. Math.} \textbf{68} (1988), 175--186.

% --- Physics Foundations ---

\bibitem{Weinberg1995}
S.~Weinberg,
\newblock \textit{The Quantum Theory of Fields, Vol.\ I},
\newblock Cambridge University Press, 1995.

\bibitem{Peskin1995}
M.~E.~Peskin and D.~V.~Schroeder,
\newblock \textit{An Introduction to Quantum Field Theory},
\newblock Addison-Wesley, 1995.

\bibitem{Penrose1971}
R.~Penrose,
\newblock ``Applications of Negative Dimensional Tensors,''
\newblock in \textit{Combinatorial Mathematics and Its Applications}, 1971.

% --- Emergence, Information, and Thermodynamics ---

\bibitem{Jacobson1995}
T.~Jacobson,
\newblock ``Thermodynamics of Spacetime: The Einstein Equation of State,''
\newblock \textit{Phys. Rev. Lett.} \textbf{75} (1995), 1260--1263.

\bibitem{Verlinde2011}
E.~Verlinde,
\newblock ``On the Origin of Gravity and the Laws of Newton,''
\newblock \textit{JHEP} \textbf{04} (2011), 029.

\bibitem{Jaynes1957}
E.~T.~Jaynes,
\newblock ``Information Theory and Statistical Mechanics,''
\newblock \textit{Phys. Rev.} \textbf{106} (1957), 620--630.

\bibitem{Shannon1948}
C.~E.~Shannon,
\newblock ``A Mathematical Theory of Communication,''
\newblock \textit{Bell System Technical Journal} \textbf{27} (1948), 379--423.

\bibitem{Landauer1961}
R.~Landauer,
\newblock ``Irreversibility and Heat Generation in the Computing Process,''
\newblock \textit{IBM J. Res. Dev.} \textbf{5} (1961), 183--191.

% --- Conceptual Foundations ---

\bibitem{Smolin2006}
L.~Smolin,
\newblock \textit{The Trouble with Physics},
\newblock Houghton Mifflin, 2006.

% --- Author's Prior Work ---

\bibitem{rsvp2024a}
Flyxion,
\newblock \textit{Constraint, Trajectory, and Irreversibility: A Unified Framework for Language, Computation, and Physics},
\newblock Independent Researcher, 2024.

\bibitem{rsvp2024b}
Flyxion,
\newblock \textit{Geometric Bayesianism with Sparse Heuristics and Thermodynamic Reachability},
\newblock Independent Researcher, 2024.

\bibitem{kes2024}
Flyxion,
\newblock \textit{The Categorical Structure of AI Alignment},
\newblock Independent Researcher, 2024.

\bibitem{spherepop2024}
Flyxion,
\newblock \textit{Identity as Event History: The Spherepop Event Calculus},
\newblock Independent Researcher, 2024.

\bibitem{rsvpcosmology2024}
Flyxion,
\newblock \textit{The RSVP Cosmological Framework: Integrating Spherepop and KES with Empirical Simulation},
\newblock Independent Researcher, 2024.

\bibitem{manifold2025}
Flyxion,
\newblock \textit{The Geometry of What Is Learned: Neural Network Semantic Manifolds},
\newblock Independent Researcher, 2025.

\bibitem{synthetic2025}
Flyxion,
\newblock \textit{Synthetic Grounding in Virtual Domains},
\newblock Independent Researcher, 2025.

\end{thebibliography}

\end{document}
